%% ma: L12-B08-0054
%% lop: 12
%% chuong: 2
%% bai: 8
%% dang: 12.8.34
%% loai: DS
%% muc: [NB, TH, VD, VD]
%% dap_an: "ĐSSS"
%% nguon: "LaTeX giáo viên gửi 21/10/2026 (M05, Vectơ và hệ trục tọa độ trong không gian, 107 câu), Phần II Câu 9"
%% kien_thuc: "Bài 8 (độ dài vectơ, hai vectơ cùng phương), lớp 10 (tính chất đường phân giác trong)"
%% trang_thai: cho-duyet
%% ly_do: "Dạng mới đề xuất 12.8.34"
%% ghi_chu: "Theo số liệu đề, AC = √86 nên b), c), d) đều Sai (rõ ràng, không mơ hồ). Nhiều khả năng tác giả định cho AC = 2AB (khi đó D(0;5/3;11/3) và b), d) Đúng); giữ nguyên số liệu đề"
\begin{ex}
Trong không gian $Oxyz$, cho tam giác $ABC$ có $A(1;2;-1)$, $B(2;-1;3)$ và $C(-4;7;5)$. Gọi $D$ là chân đường phân giác trong góc $A$ của tam giác $ABC$ ($D$ nằm trên đoạn thẳng $BC$).
\choiceTF
{\True Độ dài cạnh $AB$ bằng $\sqrt{26}$.}
{Độ dài cạnh $AC$ gấp đôi độ dài cạnh $AB$.}
{Theo tính chất đường phân giác trong, ta có hệ thức vectơ $\overrightarrow{DB}=-2\overrightarrow{DC}$.}
{Tọa độ của điểm $D$ là $\left(0;\dfrac53;\dfrac{11}3\right)$.}
\loigiai{a) $\overrightarrow{AB}=(1;-3;4)$, $AB=\sqrt{26}$. Đúng.

b) $\overrightarrow{AC}=(-5;5;6)$, $AC=\sqrt{86}\ne2\sqrt{26}$. Sai.

c) $\dfrac{DB}{DC}=\dfrac{AB}{AC}=\sqrt{\dfrac{26}{86}}\ne2$ nên $\overrightarrow{DB}\ne-2\overrightarrow{DC}$. Sai.

d) $D=\dfrac{AC\cdot B+AB\cdot C}{AB+AC}$ có tọa độ vô tỉ, trong khi $\left(0;\dfrac53;\dfrac{11}3\right)$ là điểm chia $BC$ theo tỉ số $DB:DC=1:2$. Sai.}
\end{ex}
